% ============================================================
%  Risposte alle Domande d'Esame — Corso di Meccanica
%  Professori Mazzarello e Spagnolo
% ============================================================
\documentclass[11pt,a4paper]{article}

\usepackage[utf8]{inputenc}
\usepackage[T1]{fontenc}
\usepackage[italian]{babel}
\usepackage{amsmath,amssymb,amsthm}
\usepackage{bm}
\usepackage{geometry}
\usepackage[colorlinks=true,linkcolor=black]{hyperref}
\usepackage{enumitem}
\usepackage{mathtools}
\usepackage{xcolor}

\geometry{top=2.5cm, bottom=2.5cm, left=2.5cm, right=2.5cm}

% ---- vettori e versori ----
\renewcommand{\vec}[1]{\bm{#1}}
\newcommand{\uv}[1]{\hat{\bm{#1}}}
\newcommand{\vt}{\hat{\bm{t}}}
\newcommand{\vn}{\hat{\bm{n}}}
\newcommand{\er}{\hat{\bm{e}}_r}
\newcommand{\ephi}{\hat{\bm{e}}_\varphi}
\newcommand{\ez}{\hat{\bm{e}}_z}
\newcommand{\rhat}{\hat{\bm{r}}}
\newcommand{\om}{\vec{\omega}}
\newcommand{\CM}{\mathrm{cm}}

% ---- ambienti teorema ----
\theoremstyle{plain}
\newtheorem{teorema}{Teorema}[section]
\theoremstyle{definition}
\newtheorem{definizione}{Definizione}[section]
\theoremstyle{remark}
\newtheorem{osservazione}{Osservazione}[section]

% ---- box enfasi ----
\newcommand{\importante}[1]{%
  \medskip\noindent
  \fcolorbox{blue!50!black}{blue!5}{\parbox{0.97\linewidth}{#1}}%
  \medskip}

\begin{document}

% ===============================================================
\begin{titlepage}
  \centering\vspace*{2.5cm}
  {\LARGE\bfseries Risposte alle Domande d'Esame\par}\vspace{0.6cm}
  {\Large Corso di Meccanica\par}\vspace{0.3cm}
  {\large Professori Mazzarello e Spagnolo\par}\vfill
  {\itshape Raccolta completa per argomento, con dimostrazioni e commenti fisici.}\vspace{1.2cm}
\end{titlepage}

\tableofcontents\newpage

% ===============================================================
\section{Cinematica del punto e rappresentazione intrinseca}
% ===============================================================

\subsection{Ascissa curvilinea}

Fissata una traiettoria $\mathcal{C}$ nello spazio, si sceglie un punto di
riferimento $P_0\in\mathcal{C}$ e un verso positivo di percorrenza.
\begin{definizione}
L'\textbf{ascissa curvilinea} $s(t)$ è la lunghezza d'arco con segno percorsa
dal punto $P_0$ al punto $P(t)$ lungo $\mathcal{C}$:
\[
  s(t)=\int_{P_0}^{P(t)}\mathrm{d}\ell\,.
\]
\end{definizione}
Associare a ogni punto della traiettoria uno scalare $s$ significa
\emph{parametrizzare intrisecamente} la curva: l'informazione sulla forma della
traiettoria è codificata in $\mathcal{C}$ stessa, mentre $s$ descrive
\emph{dove} si trova il punto su di essa, in modo indipendente dal sistema di
riferimento cartesiano. Il moto è quindi descritto da un'unica funzione
scalare $s(t)$ (o equivalentemente $\theta(t)$ per il pendolo, ecc.).

Per un corpo vincolato a muoversi lungo $\mathcal{C}$, lo spostamento
infinitesimo è $\mathrm{d}\vec{r}=\vt\,\mathrm{d}s$ e il lavoro elementare di
una forza $\vec{F}$ diventa:
\[
  \mathrm{d}W = \vec{F}\cdot\mathrm{d}\vec{r} = F_t\,\mathrm{d}s\,,
\]
dove $F_t=\vec{F}\cdot\vt$ è la componente tangenziale: solo questa componente
compie lavoro sul moto vincolato.

\subsection{Grandezze intrinseche: versori, raggio di curvatura e cerchio
osculatore}

\paragraph{Versore tangente.}
\[
  \vt(s)=\frac{\mathrm{d}\vec{r}}{\mathrm{d}s}\,,\qquad |\vt|=1\,.
\]
\paragraph{Velocità sempre tangente.}
Applicando la regola di derivazione composta:
\[
  \vec{v}=\frac{\mathrm{d}\vec{r}}{\mathrm{d}t}
         =\frac{\mathrm{d}\vec{r}}{\mathrm{d}s}\cdot\frac{\mathrm{d}s}{\mathrm{d}t}
         =\dot{s}\,\vt = v\,\vt\,.
\]
Poiché $\vt$ è per costruzione tangente a $\mathcal{C}$, la velocità è
\emph{sempre} tangente alla traiettoria, indipendentemente dalla sua forma.

\paragraph{Versore normale e raggio di curvatura.}
Derivando $|\vt|^2=1$ rispetto a $s$ si ottiene
$\vt\cdot\frac{\mathrm{d}\vt}{\mathrm{d}s}=0$, cioè $\frac{\mathrm{d}\vt}{\mathrm{d}s}$
è \emph{perpendicolare} a $\vt$. Si definisce il \textbf{versore normale
principale} $\vn$ e il \textbf{raggio di curvatura} $\mathcal{R}$ mediante:
\[
  \frac{\mathrm{d}\vt}{\mathrm{d}s}=\frac{\vn}{\mathcal{R}}\,,\qquad \mathcal{R}>0\,.
\]
Il \textbf{cerchio osculatore} è il cerchio di raggio $\mathcal{R}$ tangente a
$\mathcal{C}$ nel punto $P$ e con centro $C$ nella direzione di $\vn$: è il
cerchio che meglio approssima la curva al second'ordine in quel punto
(contatto di ordine 2).

\subsection{Decomposizione dell'accelerazione in componente tangenziale e
normale}

Derivando $\vec{v}=v\,\vt$ rispetto al tempo:
\[
  \vec{a}
    =\dot{v}\,\vt + v\,\dot{\vt}
    =\dot{v}\,\vt + v\,\frac{\mathrm{d}\vt}{\mathrm{d}s}\,\dot{s}
    =\underbrace{\dot{v}}_{a_t}\vt
     +\underbrace{\frac{v^2}{\mathcal{R}}}_{a_n}\vn\,.
\]
\begin{itemize}
  \item $a_t=\ddot{s}$: modifica il \emph{modulo} della velocità.
  \item $a_n=v^2/\mathcal{R}$: \textbf{accelerazione centripeta}, modifica la
    \emph{direzione} di $\vec{v}$; è diretta verso il centro del cerchio
    osculatore e compare ogni volta che la traiettoria è curva ($\mathcal{R}<\infty$),
    in qualunque sistema di riferimento inerziale.
\end{itemize}

\subsection{Derivata temporale di un vettore in sistemi fisso e mobile;
relazioni di Poisson}

Sia $S'$ un sistema di riferimento mobile che ruota con velocità angolare
$\om$ rispetto al sistema fisso $S$, e sia $\{\hat{\bm{e}}_1',\hat{\bm{e}}_2',\hat{\bm{e}}_3'\}$
la terna solidale con $S'$.

\paragraph{Perché $\dot{\hat{\bm{e}}}_i=\om\times\hat{\bm{e}}_i$.}
Poiché $|\hat{\bm{e}}_i|=1$, la derivata è perpendicolare al versore; in una
rotazione infinitesima $\mathrm{d}\vec{\varphi}=\om\,\mathrm{d}t$ lo
spostamento del versore è $\mathrm{d}\hat{\bm{e}}_i=\mathrm{d}\vec{\varphi}\times\hat{\bm{e}}_i$,
dunque $\dot{\hat{\bm{e}}}_i=\om\times\hat{\bm{e}}_i$.
Queste sono le \textbf{relazioni di Poisson}: descrivono come cambiano nel
tempo, visti dal sistema fisso, i versori solidali con il sistema mobile.

\paragraph{Derivata di un vettore generico.}
Sia $\vec{A}=\sum_i A_i'\,\hat{\bm{e}}_i'$ (componenti in $S'$). Vista da $S$:
\[
  \left.\frac{\mathrm{d}\vec{A}}{\mathrm{d}t}\right|_S
  =\left.\frac{\mathrm{d}\vec{A}}{\mathrm{d}t}\right|_{S'}+\om\times\vec{A}\,.
\]
La differenza tra le due derivate è $\om\times\vec{A}$: il termine aggiuntivo
cattura il fatto che in $S$ anche i versori di riferimento cambiano nel tempo.

\subsection{Leggi di trasformazione di velocità e accelerazione}

Sia $O'$ l'origine di $S'$, con posizione $\vec{r}_{O'}$ in $S$, e sia
$\vec{r}'$ la posizione della particella in $S'$.

\paragraph{Velocità.}
\[
  \vec{v}=\dot{\vec{r}}_{O'}+\left.\frac{\mathrm{d}\vec{r}'}{\mathrm{d}t}\right|_S
         =\vec{v}_{O'}+\vec{v}'+\om\times\vec{r}'
         \equiv \vec{v}'+\vec{v}_\mathrm{tr}\,,
\]
dove $\vec{v}_\mathrm{tr}=\vec{v}_{O'}+\om\times\vec{r}'$ è la
\textbf{velocità di trascinamento} (velocità che avrebbe un punto solidale con
$S'$ che istantaneamente coincide con la particella).

\paragraph{Accelerazione.}
Derivando ancora:
\[
  \vec{a}=\vec{a}'+\underbrace{\dot{\vec{v}}_{O'}
          +\dot{\om}\times\vec{r}'+\om\times(\om\times\vec{r}')}_{\vec{a}_\mathrm{tr}}
          +2\om\times\vec{v}'\,.
\]
Il termine $2\om\times\vec{v}'$ è l'\textbf{accelerazione di Coriolis};
$\vec{a}_\mathrm{tr}$ è l'accelerazione di trascinamento.

\subsection{Coordinate cilindriche: versori e loro derivate}

In coordinate $(r,\varphi,z)$, la terna ortonormale mobile è $\{\er,\ephi,\ez\}$
con $\er$ radiale, $\ephi$ azimutale, $\ez$ assiale. Le derivate temporali,
usando le relazioni di Poisson con $\om=\dot\varphi\,\ez$, sono:
\[
  \dot{\er}=\dot\varphi\,\ephi\,,\qquad
  \dot{\ephi}=-\dot\varphi\,\er\,,\qquad
  \dot{\ez}=\vec{0}\,.
\]
Derivando $\vec{r}=r\,\er+z\,\ez$:
\begin{align*}
  \vec{v}&=\dot r\,\er + r\dot\varphi\,\ephi + \dot z\,\ez\,,\\
  \vec{a}&=(\ddot r-r\dot\varphi^2)\,\er+(r\ddot\varphi+2\dot r\dot\varphi)\,\ephi+\ddot z\,\ez\,.
\end{align*}
Il termine $-r\dot\varphi^2\,\er$ è l'accelerazione centripeta; il termine
$2\dot r\dot\varphi\,\ephi$ è l'accelerazione di Coriolis intrinseca delle
coordinate cilindriche (non forza, ma termine cinematico).

% ===============================================================
\section{Sistemi di riferimento non inerziali e forze apparenti}
% ===============================================================

\subsection{Sistemi inerziali e non inerziali}

Un sistema di riferimento è \textbf{inerziale} se in esso i principi della
dinamica di Newton valgono nella loro forma originale (in particolare, un
corpo libero si muove di moto rettilineo uniforme). Un sistema $S'$ in moto
traslatorio accelerato o in rotazione rispetto a un sistema inerziale $S$ è
\textbf{non inerziale}.

\subsection{Seconda legge in un sistema non inerziale e forze apparenti}

Dalla legge di trasformazione dell'accelerazione, sostituendo in $\vec{F}=m\vec{a}$
(valida nel sistema inerziale $S$):
\[
  m\vec{a}'=\vec{F}
            -m\vec{a}_{O'}
            -m\dot{\om}\times\vec{r}'
            -m\om\times(\om\times\vec{r}')
            -2m\om\times\vec{v}'\,.
\]
Le quantità aggiuntive, che \emph{non} derivano da interazioni fisiche reali
ma dalla non inerzia del sistema, sono le \textbf{forze apparenti} (pseudo-forze):
\begin{itemize}
  \item $-m\vec{a}_{O'}$: forza di trascinamento traslazionale.
  \item $-m\dot{\om}\times\vec{r}'$: forza di Eulero (compare se $|\om|$ o la
    sua direzione variano nel tempo).
  \item $\vec{F}_\mathrm{cf}=-m\om\times(\om\times\vec{r}')$: \textbf{forza centrifuga}.
  \item $\vec{F}_\mathrm{Cor}=-2m\om\times\vec{v}'$: \textbf{forza di Coriolis}.
\end{itemize}

\subsection{Accelerazione e velocità di trascinamento}

L'\textbf{accelerazione di trascinamento}
$\vec{a}_\mathrm{tr}=\vec{a}_{O'}+\dot{\om}\times\vec{r}'+\om\times(\om\times\vec{r}')$
è l'accelerazione che avrebbe un punto \emph{solidale} con $S'$, coincidente
istantaneamente con la particella. Essa tiene conto della traslazione
accelerata di $O'$, della variazione di $\om$ e dell'accelerazione centripeta
dovuta alla rotazione.

\subsection{Forza centrifuga: definizione, significato fisico e conservatività}

\[
  \vec{F}_\mathrm{cf}=-m\om\times(\om\times\vec{r}')=m\omega^2\vec{\rho}'\,,
\]
dove $\vec{\rho}'$ è la componente di $\vec{r}'$ perpendicolare all'asse di
rotazione, diretta verso l'esterno. La forza centrifuga è quindi radiale
rispetto all'asse, diretta centrifugamente, e aumenta con la distanza dall'asse.

\paragraph{Conservatività.}
Con l'asse di rotazione orientato lungo $z$, si ha
$\vec{F}_\mathrm{cf}=m\omega^2(x'\hat{\bm{x}}+y'\hat{\bm{y}})
=-\nabla\left(-\tfrac{1}{2}m\omega^2\rho'^2\right)$,
dunque la forza centrifuga è conservativa con \textbf{potenziale centrifugo}:
\[
  U_\mathrm{cf}=-\frac{1}{2}m\omega^2\rho'^2\,.
\]
La condizione necessaria è che $\om$ sia costante (modulo e direzione). Se
$\om$ varia nel tempo, compare la forza di Eulero che, in generale, non è
conservativa.

\subsection{Forza di Coriolis}

\[
  \vec{F}_\mathrm{Cor}=-2m\om\times\vec{v}'\,.
\]
È perpendicolare sia alla velocità relativa $\vec{v}'$ sia all'asse $\om$.
\paragraph{La forza di Coriolis non compie lavoro:}
\[
  P_\mathrm{Cor}=\vec{F}_\mathrm{Cor}\cdot\vec{v}'
               =-2m(\om\times\vec{v}')\cdot\vec{v}'=0\,,
\]
perché il prodotto misto $(\vec{A}\times\vec{B})\cdot\vec{B}=0$ sempre.
La forza di Coriolis devia la traiettoria ma non modifica la velocità scalare.
Fisicamente è responsabile della deviazione dei venti (cicloni, anticicloni)
e della rotazione del piano di Foucault.

\subsection{Forze apparenti in coordinate cilindriche}

Per un sistema ruotante con $\om=\omega\ez$, le componenti delle forze
apparenti in coordinate cilindriche sono:
\begin{align*}
  \text{centrifuga: }&\vec{F}_\mathrm{cf}=m\omega^2 r'\,\er\,,\\
  \text{Coriolis: }&\vec{F}_\mathrm{Cor}=-2m\omega\dot r'\,\ephi + 2m\omega r'\dot\varphi'\,\er\,.
\end{align*}

% ===============================================================
\section{Forze, lavoro, conservatività ed energia potenziale}
% ===============================================================

\subsection{Lavoro e lavoro infinitesimo}

Il \textbf{lavoro} compiuto da $\vec{F}$ lungo un percorso da $A$ a $B$:
\[
  W_{A\to B}=\int_A^B\vec{F}\cdot\mathrm{d}\vec{r}\,.
\]
Il \textbf{lavoro infinitesimo} $\mathrm{d}W=\vec{F}\cdot\mathrm{d}\vec{r}$
può essere riscritto in coordinate cartesiane come
$F_x\,\mathrm{d}x+F_y\,\mathrm{d}y+F_z\,\mathrm{d}z$.

\subsection{Definizione e caratteristiche delle forze conservative}

\begin{definizione}
Una forza $\vec{F}$ è \textbf{conservativa} se il lavoro lungo qualunque
percorso chiuso è nullo:
$\oint_\gamma\vec{F}\cdot\mathrm{d}\vec{r}=0$ per ogni curva chiusa $\gamma$.
\end{definizione}
Le seguenti condizioni sono equivalenti:
\begin{enumerate}[label=(\roman*)]
  \item $\oint\vec{F}\cdot\mathrm{d}\vec{r}=0$ (circuitazione nulla).
  \item $\nabla\times\vec{F}=\vec{0}$ (rotore nullo).
  \item $\exists\, U(\vec{r}):\;\vec{F}=-\nabla U$ (esiste un potenziale).
\end{enumerate}

\paragraph{Da circuitazione nulla a rotore nullo.}
Per il teorema di Stokes, per qualsiasi superficie $\Sigma$ con bordo $\partial\Sigma$:
\[
  \oint_{\partial\Sigma}\vec{F}\cdot\mathrm{d}\vec{r}
  =\iint_\Sigma(\nabla\times\vec{F})\cdot\mathrm{d}\vec{S}\,.
\]
Se la circuitazione è nulla per ogni curva chiusa $\partial\Sigma$, l'integrale
del rotore è nullo per ogni superficie $\Sigma$: poiché $\Sigma$ è arbitraria,
si conclude $\nabla\times\vec{F}=\vec{0}$ in ogni punto.

\paragraph{Forza conservativa come gradiente di un potenziale.}
Dal rotore nullo (in un dominio semplicemente connesso) segue che
$\vec{F}\cdot\mathrm{d}\vec{r}$ è una differenziale esatta, e si può definire
$U(\vec{r})=-\int_{\vec{r}_0}^{\vec{r}}\vec{F}\cdot\mathrm{d}\vec{r}'$,
da cui $\vec{F}=-\nabla U$.

\subsection{Dimostrazione che forze specifiche sono (o non sono) conservative}

\paragraph{Forza centrale a simmetria sferica: $\vec{F}=f(r)\rhat$.}
\[
  \mathrm{d}W=\vec{F}\cdot\mathrm{d}\vec{r}=f(r)\rhat\cdot\mathrm{d}\vec{r}=f(r)\,\mathrm{d}r\,,
\]
che è differenziale esatta di $U(r)=-\int f(r)\,\mathrm{d}r$. Dunque è
conservativa. Equivalentemente, $\nabla\times(f(r)\rhat)=0$ per simmetria
sferica (verifica diretta in coordinate sferiche).

\paragraph{Forza centrale a simmetria cilindrica: $\vec{F}=f(\rho)\er$.}
Analogamente, $\mathrm{d}W=f(\rho)\,\mathrm{d}\rho$ è esatta, con
$U=-\int f(\rho)\,\mathrm{d}\rho$.

\paragraph{Forza elastica: $\vec{F}=-kx\,\hat{\bm{x}}$.}
$\nabla\times\vec{F}=\vec{0}$ (componenti non dipendono dalle variabili
perpendicolari all'asse), con $U=\tfrac{1}{2}kx^2$.

\paragraph{Attrito viscoso: $\vec{F}=-b\vec{v}$.}
La forza dipende dalla velocità, non dalla posizione: non è definita nello
spazio delle posizioni, dunque non può essere scritta come gradiente di un
potenziale. Il rotore non è nullo (la forza cambia direzione con $\vec{v}$).
Il lavoro lungo un ciclo non è nullo e dipende dal percorso specifico.

\paragraph{Attrito radente dinamico.}
$\vec{F}_d=-\mu_d N\,\hat{v}$: fa sempre lavoro negativo ($\mathrm{d}W<0$),
dissipa energia cinetica.

\subsection{Energia potenziale, equilibrio stabile e instabile}

\begin{definizione}
L'\textbf{energia potenziale} $U(\vec{r})$ è la funzione scalare definita (a
meno di una costante) da $\vec{F}=-\nabla U$.
\end{definizione}

Una posizione $x_0$ è di \textbf{equilibrio} se $F(x_0)=0$, cioè
$U'(x_0)=0$.

\paragraph{Studio della stabilità via espansione di Taylor.}
Intorno a $x_0$, con $\xi=x-x_0$:
\[
  U(x_0+\xi)=U(x_0)+\frac{1}{2}U''(x_0)\xi^2+O(\xi^3)\,,\qquad
  F(\xi)\approx -U''(x_0)\xi\,.
\]

\begin{itemize}
  \item \textbf{Equilibrio stabile} se $U''(x_0)>0$ (minimo di $U$): la forza
    è di richiamo ($F=-k_\mathrm{eff}\xi$, $k_\mathrm{eff}>0$). Una perturbazione
    piccola genera moto oscillatorio attorno a $x_0$.
  \item \textbf{Equilibrio instabile} se $U''(x_0)<0$ (massimo di $U$): la forza
    è repulsiva, allontana il punto da $x_0$. Una perturbazione piccola cresce
    esponenzialmente.
\end{itemize}

\paragraph{Perché la derivata seconda determina la stabilità.}
Fisicamente: attorno a un minimo di $U$, qualunque spostamento dal punto di
equilibrio aumenta $U$ e la forza conservativa richiama il sistema verso il
minimo. Attorno a un massimo, qualunque spostamento diminuisce $U$ e la forza
spinge il sistema sempre più lontano. Matematicamente, il segno di $U''$
determina se la costante di richiamo effettiva è positiva (oscillazione) o
negativa (divergenza).

\subsection{Moto intorno ai punti di equilibrio}

Attorno a un equilibrio stabile, l'equazione del moto è
$m\ddot\xi=-U''(x_0)\xi$, cioè quella di un oscillatore armonico con
$\omega_0=\sqrt{U''(x_0)/m}$. Attorno a un equilibrio instabile,
$m\ddot\xi=|U''(x_0)|\xi$ con soluzione $\xi(t)=A e^{\alpha t}+Be^{-\alpha t}$,
$\alpha=\sqrt{|U''(x_0)|/m}$.

\subsection{Teorema delle forze vive (lavoro–energia)}

\begin{teorema}[Forze vive]
Il lavoro totale compiuto su un punto materiale è uguale alla variazione
della sua energia cinetica:
\[
  W_{A\to B}=T_B-T_A=\Delta T\,.
\]
\end{teorema}
\begin{proof}
Da $m\ddot{\vec{r}}=\vec{F}_\mathrm{tot}$, si moltiplica scalarmente per
$\vec{v}$:
\[
  \vec{F}_\mathrm{tot}\cdot\vec{v}=m\vec{v}\cdot\dot{\vec{v}}=\frac{\mathrm{d}}{\mathrm{d}t}\!\left(\frac{mv^2}{2}\right)=\frac{\mathrm{d}T}{\mathrm{d}t}\,.
\]
Integrando nel tempo da $A$ a $B$:
$W_{A\to B}=\int_{t_A}^{t_B}\vec{F}\cdot\vec{v}\,\mathrm{d}t=T_B-T_A$.
\end{proof}

\paragraph{Conservazione dell'energia meccanica.}
Per forze conservative $\vec{F}=-\nabla U$:
$W_{A\to B}=U_A-U_B$, dunque $T_B-T_A=U_A-U_B$, ovvero:
\[
  E=T+U=\mathrm{const}\,.
\]

% ===============================================================
\section{Oscillazioni e pendolo semplice}
% ===============================================================

\subsection{Oscillatore armonico: equazione del moto e significato dei
parametri}

L'equazione dell'oscillatore armonico semplice è:
\[
  m\ddot x+kx=0 \implies \ddot x+\omega_0^2 x=0\,,\quad\omega_0=\sqrt{\frac{k}{m}}\,.
\]
La soluzione generale è:
\[
  x(t)=A\cos(\omega_0 t+\phi)\,,
\]
con \textbf{ampiezza} $A$ e \textbf{fase iniziale} $\phi$ determinati dalle
condizioni iniziali. Il parametro $\omega_0$ è la \textbf{pulsazione propria}:
il moto è periodico con periodo $T_0=2\pi/\omega_0$ e frequenza
$\nu_0=\omega_0/(2\pi)$.

\paragraph{Come si arriva all'oscillatore armonico da un equilibrio stabile.}
Attorno a un equilibrio stabile $x_0$ (con $U'(x_0)=0$, $U''(x_0)>0$),
ponendo $\xi=x-x_0$ e sviluppando al second'ordine:
\[
  U(\xi)\approx U_0+\tfrac{1}{2}k_\mathrm{eff}\xi^2\,,\quad
  F(\xi)\approx -k_\mathrm{eff}\xi\,,\quad
  k_\mathrm{eff}=U''(x_0)>0\,.
\]
L'equazione del moto $m\ddot\xi=-k_\mathrm{eff}\xi$ è quella dell'oscillatore
armonico con $\omega_0=\sqrt{k_\mathrm{eff}/m}$.

\paragraph{Equilibrio instabile.}
Con $U''(x_0)<0$, la forza è $F=|k_\mathrm{eff}|\xi$, con soluzione
$\xi=Ae^{\alpha t}+Be^{-\alpha t}$ ($\alpha>0$): moto esponenzialmente
divergente, non oscillatorio.

\subsection{Oscillatore smorzato: regimi e attraversamento dell'equilibrio}

Con smorzamento viscoso $-b\dot x$:
\[
  \ddot x+2\gamma\dot x+\omega_0^2 x=0\,,\qquad\gamma=\frac{b}{2m}\,.
\]
L'equazione caratteristica è $\lambda^2+2\gamma\lambda+\omega_0^2=0$, con
$\Delta=\gamma^2-\omega_0^2$.

\begin{description}[leftmargin=2em]
  \item[Sottosmorzato ($\gamma<\omega_0$, $\Delta<0$):]
    $\lambda=-\gamma\pm i\omega_1$, $\omega_1=\sqrt{\omega_0^2-\gamma^2}$.
    \[x(t)=e^{-\gamma t}(A\cos\omega_1 t+B\sin\omega_1 t)\,.\]
    Moto oscillatorio con ampiezza decrescente; il sistema \emph{passa infinite
    volte} per l'equilibrio.

  \item[Criticamente smorzato ($\gamma=\omega_0$, $\Delta=0$):]
    $\lambda=-\gamma$ radice doppia.
    \[x(t)=(A+Bt)e^{-\gamma t}\,.\]
    Nessuna oscillazione; il sistema raggiunge l'equilibrio nel \emph{tempo
    minimo}. Può passare per $x=0$ \emph{al più una volta}, e solo se $A$ e
    $B$ hanno segni opportuni.

  \item[Sovrasmorzato ($\gamma>\omega_0$, $\Delta>0$):]
    $\lambda_{1,2}=-\gamma\pm\beta$, $\beta=\sqrt{\gamma^2-\omega_0^2}>0$.
    \[x(t)=e^{-\gamma t}(Ae^{\beta t}+Be^{-\beta t})\,.\]
    Nessuna oscillazione; il ritorno all'equilibrio è più lento che nel caso
    critico. Può passare per $x=0$ \emph{al più una volta} o \emph{mai},
    secondo le condizioni iniziali.
\end{description}

\subsection{Oscillatore forzato e risonanza}

Con forzante armonica:
\[
  \ddot x+2\gamma\dot x+\omega_0^2 x=\frac{F_0}{m}\cos(\omega t)\,.
\]
In regime stazionario (transitorio esaurito) la soluzione particolare è
$x_p=D\cos(\omega t+\delta)$ con:
\[
  D(\omega)=\frac{F_0/m}{\sqrt{(\omega_0^2-\omega^2)^2+4\gamma^2\omega^2}}\,.
\]
$D(\omega)$ è massima (condizione di \textbf{risonanza}) a:
\[
  \omega_r=\sqrt{\omega_0^2-2\gamma^2}<\omega_0\quad
  (\text{per }\gamma<\omega_0/\sqrt{2})\,.
\]
Per $\gamma\to 0$ l'ampiezza di risonanza diverge: $D\to\infty$.

\subsection{Pendolo semplice}

Massa puntiforme $m$ appesa a filo ideale di lunghezza $L$, angolo $\theta$
dalla verticale.

\paragraph{Via rappresentazione intrinseca.}
L'ascissa curvilinea è $s=L\theta$. La componente tangenziale del peso è
$F_t=-mg\sin\theta$. La componente normale è bilanciata dalla tensione.
\[
  m\ddot s=-mg\sin\theta \implies L\ddot\theta=-g\sin\theta
  \implies \ddot\theta+\frac{g}{L}\sin\theta=0\,.
\]

\paragraph{Via momento angolare rispetto al fulcro $O$.}
$\vec{L}_O=mL^2\dot\theta\,\hat{\bm{e}}_\perp$;
$\vec{\tau}_O=-mgL\sin\theta\,\hat{\bm{e}}_\perp$.
La seconda equazione cardinale $\dot{\vec{L}}_O=\vec{\tau}_O$ dà:
$mL^2\ddot\theta=-mgL\sin\theta$, ovvero la stessa equazione.

\paragraph{Via energia meccanica.}
$E=\tfrac{1}{2}mL^2\dot\theta^2+mgL(1-\cos\theta)=\mathrm{const}$.
Derivando rispetto al tempo:
$mL^2\dot\theta\ddot\theta+mgL\sin\theta\,\dot\theta=0\implies\ddot\theta=-\tfrac{g}{L}\sin\theta$.

\paragraph{Approssimazione per piccoli angoli.}
Per $|\theta|\ll1$: $\sin\theta\approx\theta$, $\cos\theta\approx1-\theta^2/2$.
L'equazione diventa lineare:
\[
  \ddot\theta+\frac{g}{L}\theta=0\,,\qquad
  \omega_0=\sqrt{\frac{g}{L}}\,,\qquad
  T_0=2\pi\sqrt{\frac{L}{g}}\,.
\]
L'approssimazione è necessaria perché l'equazione originale $\ddot\theta=-\tfrac{g}{L}\sin\theta$
è \emph{non lineare} e non ammette soluzioni analitiche elementari; la
linearizzazione la riconduce all'oscillatore armonico con periodo
indipendente dall'ampiezza (isocronia di Galileo), valida per
$|\theta|\lesssim 10^\circ$.

% ===============================================================
\section{Sistemi di punti, centro di massa ed equazioni cardinali}
% ===============================================================

\subsection{Centro di massa: definizione e proprietà cinematiche}

\begin{definizione}
Il \textbf{centro di massa} di un sistema di $N$ punti materiali è:
\[
  \vec{r}_\CM=\frac{1}{M}\sum_{i=1}^N m_i\vec{r}_i\,,\qquad M=\sum_i m_i\,.
\]
Per un corpo continuo: $\vec{r}_\CM=\frac{1}{M}\int\vec{r}\,\mathrm{d}m$.
\end{definizione}

La \textbf{quantità di moto totale} è:
\[
  \vec{Q}=\sum_i m_i\vec{v}_i=M\dot{\vec{r}}_\CM=M\vec{v}_\CM\,.
\]
Il CM si muove come se tutta la massa $M$ fosse concentrata in esso.

\paragraph{Centro di massa con asse di simmetria.}
Se il sistema ha un asse di simmetria geometrica, le contribuzioni di masse
simmetricamente disposte rispetto all'asse si cancellano nelle direzioni
perpendicolari, quindi il CM giace sull'asse di simmetria.

\subsection{Prima equazione cardinale}

\begin{teorema}[Prima equazione cardinale]
\[
  \dot{\vec{Q}}=M\vec{a}_\CM=\vec{F}^\mathrm{est}\,.
\]
\end{teorema}
\begin{proof}
Per il secondo principio applicato a ciascun punto:
$m_i\vec{a}_i=\vec{F}_i^\mathrm{est}+\sum_{j\neq i}\vec{F}_{ij}$.
Sommando su tutti i punti:
$M\vec{a}_\CM=\vec{F}^\mathrm{est}+\sum_{i}\sum_{j\neq i}\vec{F}_{ij}$.
Per il terzo principio di Newton $\vec{F}_{ij}+\vec{F}_{ji}=\vec{0}$, quindi
$\sum_{i,j,i\neq j}\vec{F}_{ij}=\vec{0}$ e le \emph{forze interne si elidono}.
\end{proof}
Il CM si muove come un punto materiale di massa $M$ soggetto alla sola
risultante delle forze esterne. Un sistema può ruotare senza risultante
esterna nulla, ma il suo CM mantiene il moto rettilineo uniforme (o stazionario).

\subsection{Seconda equazione cardinale}

\[
  \dot{\vec{L}}_O=\vec{M}_O^\mathrm{est}\,,
\]
dove $\vec{L}_O=\sum_i\vec{r}_i\times m_i\vec{v}_i$ e $\vec{M}_O^\mathrm{est}$
è la risultante dei momenti delle forze esterne rispetto a $O$. Le forze
interne, se centrali (terzo principio in forma forte), non contribuiscono.

\subsection{Teoremi di König}

Sia $\vec{v}_i=\vec{v}_\CM+\vec{v}'_i$ la scomposizione tra moto del CM e
moto relativo al CM (per definizione $\sum_i m_i\vec{v}'_i=\vec{0}$ e
$\sum_i m_i\vec{r}'_i=\vec{0}$).

\begin{teorema}[König per l'energia cinetica]
\[
  T=\underbrace{\frac{1}{2}Mv_\CM^2}_{T_\CM}+\underbrace{\sum_i\frac{1}{2}m_iv_i'^2}_{T'}\,.
\]
\end{teorema}
\begin{proof}
$T=\tfrac{1}{2}\sum_i m_i|\vec{v}_\CM+\vec{v}'_i|^2
  =\tfrac{1}{2}Mv_\CM^2
  +\vec{v}_\CM\cdot\underbrace{\sum_i m_i\vec{v}'_i}_{=\vec{0}}
  +T'$.
\end{proof}
L'energia cinetica totale si separa nel contributo del moto del CM e nel
contributo del moto relativo al CM (energia cinetica interna).

\begin{teorema}[König per il momento angolare]
\[
  \vec{L}_O=\underbrace{\vec{r}_\CM\times M\vec{v}_\CM}_{\text{orbitale}}
            +\underbrace{\sum_i m_i\vec{r}'_i\times\vec{v}'_i}_{\vec{L}'=\text{spin}}\,.
\]
\end{teorema}
\begin{proof}
$\vec{L}_O=\sum_i m_i(\vec{r}_\CM+\vec{r}'_i)\times(\vec{v}_\CM+\vec{v}'_i)$.
I termini misti $\vec{r}_\CM\times\sum m_i\vec{v}'_i$ e
$(\sum m_i\vec{r}'_i)\times\vec{v}_\CM$ sono nulli per la proprietà del CM.
\end{proof}

% ===============================================================
\section{Corpo rigido e rotazione attorno ad asse fisso}
% ===============================================================

\subsection{Cinematica del moto rigido rispetto ad asse fisso}

\begin{definizione}
Un \textbf{corpo rigido} è un sistema di punti per cui
$|\vec{r}_i-\vec{r}_j|=\mathrm{const}$ per ogni coppia $(i,j)$.
\end{definizione}

Per la rotazione attorno all'asse fisso $\hat{\bm{z}}$, con velocità angolare
$\om=\omega\ez$:
\[
  \vec{v}_i=\om\times\vec{r}_i\,,\qquad |\vec{v}_i|=\omega\rho_i\,,
\]
dove $\rho_i$ è la distanza del punto $i$ dall'asse. L'equazione del moto di
ogni punto del corpo può essere scritta come
$\vec{r}_i(t)=R_z(\theta(t))\,\vec{r}_i(0)$, con $\dot\theta=\omega$.

\paragraph{Quando può ruotare senza forze esterne?}
La prima equazione cardinale richiede $\vec{F}^\mathrm{est}=M\vec{a}_\CM$.
Se il CM si trova sull'asse di rotazione, $\vec{a}_\CM=\vec{0}$ e non è
necessaria alcuna forza esterna netta per mantenere il moto. Se il CM è
fuori dall'asse, occorre una forza centripeta che lo mantenga in rotazione.

\paragraph{Quando può ruotare senza momenti esterni?}
La seconda equazione cardinale richiede $\vec{M}^\mathrm{est}=\dot{\vec{L}}$.
Se l'asse è un asse principale di inerzia (vedi oltre), $\vec{L}$ è parallelo
a $\vec{\omega}$, $\dot{\vec{L}}$ è nullo se $\omega=\mathrm{const}$, e non
occorre alcun momento esterno.

\subsection{Momento di inerzia: definizione e significato fisico}

\begin{definizione}
Il \textbf{momento di inerzia} rispetto all'asse $z$ è:
\[
  I=\sum_i m_i\rho_i^2=\int\rho^2\,\mathrm{d}m\,.
\]
\end{definizione}
Misura la resistenza del corpo alla variazione della velocità angolare: è
l'analogo rotazionale della massa. L'energia cinetica di rotazione è
$T_\mathrm{rot}=\tfrac{1}{2}I\omega^2$.

\subsection{Momento angolare parallelo e perpendicolare all'asse}

Il momento angolare totale rispetto a un polo $P$ sull'asse di rotazione è:
\[
  \vec{L}_P=\sum_i m_i\vec{r}_{Pi}\times(\om\times\vec{r}_{Pi})
           =\sum_i m_i\bigl[\om\,r_{Pi}^2-\vec{r}_{Pi}(\vec{r}_{Pi}\cdot\om)\bigr]\,.
\]
Scomponendo $\vec{r}_{Pi}=\vec{\rho}_i+z_i\ez$ ($\vec{\rho}_i$ perpendicolare
all'asse):

\begin{itemize}
  \item Componente \textbf{parallela} all'asse:
    \[L_z=\sum_i m_i\rho_i^2\,\omega=I\omega\,.\]
  \item Componente \textbf{perpendicolare} all'asse:
    \[
      \vec{L}_\perp=-\omega\sum_i m_i z_i\vec{\rho}_i\,.
    \]
\end{itemize}

\paragraph{Quando si annulla la componente perpendicolare?}
$\vec{L}_\perp=\vec{0}$ sse $\sum_i m_i z_i\vec{\rho}_i=\vec{0}$, cioè
quando l'asse di rotazione è un \textbf{asse principale di inerzia} (autovettore
del tensore di inerzia $\overline{\overline{I}}$, con componenti
$I_{\alpha\beta}=\sum_i m_i(\delta_{\alpha\beta}r_i^2-r_{i\alpha}r_{i\beta})$).
Gli assi di simmetria geometrica sono automaticamente assi principali.

Se $\vec{L}_\perp\neq\vec{0}$, il vettore $\vec{L}$ precede attorno all'asse
e i cuscinetti devono esercitare una coppia per mantenere l'asse fisso.

\paragraph{CM sull'asse di rotazione.}
Se il CM è sull'asse, la forza centripeta necessaria a mantenerlo in orbita è
nulla: la reazione traslazionale dei cuscinetti è nulla. Se in più l'asse è
principale, anche la coppia dei cuscinetti è nulla: il corpo ruota senza
alcun vincolo attivo.

\subsection{Seconda equazione cardinale per la rotazione}

Rispetto al CM (o a un polo fisso sull'asse):
\[
  \dot{L}_z=\tau_z^\mathrm{est} \implies I\dot\omega=\tau_z^\mathrm{est}\,.
\]
Questa equazione scalare è l'analogo rotazionale di $F=ma$.

\subsection{Teorema di Poinsot}

\begin{teorema}[Poinsot]
Un corpo rigido libero (nessun momento esterno) evolve in modo che il vettore
$\vec{L}$ rimanga costante nello spazio, mentre $\vec{\omega}$ descrive un
cono nello spazio (\emph{cono spaziale}) e un cono nel corpo (\emph{cono del
corpo}), che rotolano l'uno sull'altro senza scorrimento.
\end{teorema}

Poiché $\dot{\vec{L}}=\vec{0}$, il momento angolare è fisso. Tuttavia
$\vec{\omega}$ non è in generale parallelo a $\vec{L}$ (solo per rotazione
attorno a un asse principale): questo genera i \textbf{moti giroscopici}
(precessione e nutazione) osservati in trottole libere e giroscopi.

% ===============================================================
\section{Puro rotolamento, attrito statico e teorema di Huygens-Steiner}
% ===============================================================

\subsection{Moto di puro rotolamento}

\begin{definizione}
Un corpo di raggio $R$ è in \textbf{puro rotolamento} su una superficie piana
se il punto di contatto ha velocità istantanea nulla:
\[
  v_\CM=\omega R\,.
\]
\end{definizione}

\paragraph{Dimostrazione che il punto di contatto $C$ è ferma.}
$\vec{v}_C=\vec{v}_\CM+\om\times(\vec{r}_C-\vec{r}_\CM)=v_\CM\hat{\bm{x}}+\omega\ez\times(-R\hat{\bm{y}})=v_\CM\hat{\bm{x}}-\omega R\hat{\bm{x}}=\vec{0}$ (usando $v_\CM=\omega R$).

\paragraph{Asse di istantanea rotazione.}
Poiché $C$ è istantaneamente fermo, il corpo ruota istantaneamente attorno
all'asse passante per $C$: questo è l'\textbf{asse di istantanea rotazione}.
La velocità di un qualunque punto $P$ del corpo vale:
$\vec{v}_P=\om\times(\vec{r}_P-\vec{r}_C)$, come per una pura rotazione
attorno a $C$.

\subsection{Energia cinetica nel puro rotolamento}

Dal teorema di König:
\[
  T=\frac{1}{2}Mv_\CM^2+\frac{1}{2}I_\CM\omega^2
   =\frac{1}{2}\!\left(M+\frac{I_\CM}{R^2}\right)v_\CM^2\,.
\]
Usando l'asse di istantanea rotazione e il teorema di Huygens-Steiner
($I_C=I_\CM+MR^2$):
\[
  T=\frac{1}{2}I_C\omega^2\,.
\]

\subsection{Equazioni cardinali e ruolo dell'attrito}

La prima equazione cardinale (asse $x$):
\[
  Ma_\CM=F^\mathrm{est}-f\,,
\]
la seconda rispetto al CM:
\[
  I_\CM\alpha=fR\,,
\]
con il vincolo $a_\CM=\alpha R$.

L'attrito statico $f$ è \emph{essenziale} per il puro rotolamento: fornisce
la coppia che produce l'accelerazione angolare. La componente $\phi_x$ della
reazione vincolare (rispetto all'asse orizzontale) corrisponde proprio
all'attrito statico. L'attrito non compie lavoro perché il punto di contatto
è istantaneamente fermo.

\paragraph{Convenienza del polo $C$.}
Calcolando i momenti rispetto al punto di contatto $C$, l'attrito ha braccio
nullo e scompare dall'equazione:
$I_C\alpha=\tau_C^\text{(forze esterne senza attrito)}$.

\subsection{Teorema di Huygens-Steiner}

\begin{teorema}[Huygens-Steiner o degli assi paralleli]
Dati due assi paralleli, uno passante per il CM a distanza $d$ dall'altro:
\[
  I_P=I_\CM+Md^2\,.
\]
\end{teorema}
\begin{proof}
Sia $\vec{\rho}'_i=\vec{\rho}_i-\vec{\rho}_\CM$ la proiezione perpendicolare
all'asse di ogni punto nel sistema del CM, e $\vec{d}=\vec{\rho}_\CM-\vec{\rho}_P$.
\begin{align*}
  I_P&=\sum_i m_i|\vec{\rho}_i-\vec{\rho}_P|^2
      =\sum_i m_i|\vec{\rho}'_i+\vec{d}|^2\\
     &=\sum_i m_i\rho_i'^2
       +2\vec{d}\cdot\underbrace{\sum_i m_i\vec{\rho}'_i}_{=\vec{0}}
       +Md^2
      =I_\CM+Md^2\,.
\end{align*}
Il termine lineare si annulla perché $\sum_i m_i\vec{\rho}'_i=\vec{0}$ per
definizione del CM.
\end{proof}
Conseguenza: fra tutti gli assi paralleli a una data direzione, $I$ è minima
per l'asse passante per il CM.

% ===============================================================
\section{Gravitazione, campi centrali e leggi di Keplero}
% ===============================================================

\subsection{Forza gravitazionale come forza centrale conservativa}

\[
  \vec{F}=-\frac{GMm}{r^2}\rhat\,,
\]
è una forza centrale a simmetria sferica, quindi conservativa (cfr.\
Sez.~\ref{sec:forze}) con potenziale:
\[
  U(r)=-\frac{GMm}{r}\,.
\]

\subsection{Conservazione del momento angolare nei campi centrali}

Per qualunque forza centrale $\vec{F}=f(r)\rhat$:
\[
  \dot{\vec{L}}=\vec{r}\times\vec{F}=\vec{r}\times f(r)\rhat=\vec{0}\implies\vec{L}=\mathrm{const}\,.
\]
Il moto è quindi confinato nel piano perpendicolare a $\vec{L}$.

\subsection{Energia meccanica e potenziale efficace}

In coordinate polari nel piano orbitale:
\[
  E=\frac{1}{2}m\dot r^2+\frac{L^2}{2mr^2}-\frac{GMm}{r}
   =\frac{1}{2}m\dot r^2+U_\mathrm{eff}(r)\,.
\]
Il \textbf{potenziale efficace} $U_\mathrm{eff}=\frac{L^2}{2mr^2}-\frac{GMm}{r}$
ha un minimo che corrisponde all'orbita circolare; orbite legate ($E<0$)
sono ellittiche.

\subsection{Leggi di Keplero}

\begin{enumerate}
  \item \textbf{Prima legge (orbite ellittiche):} Ogni pianeta orbita su
    un'ellisse con il Sole in uno dei fuochi. Segue dall'integrazione
    dell'equazione del moto con $F\propto1/r^2$.

  \item \textbf{Seconda legge (legge delle aree):} Il raggio vettore
    pianeta-Sole spazza aree uguali in tempi uguali:
    \[
      \frac{\mathrm{d}A}{\mathrm{d}t}=\frac{L}{2m}=\mathrm{const}\,.
    \]
    È la diretta conseguenza della conservazione del momento angolare:
    $\mathrm{d}A=\tfrac{1}{2}|r\,\mathrm{d}\varphi|\cdot r=\tfrac{1}{2}r^2\,\mathrm{d}\varphi$
    e $L=mr^2\dot\varphi$.

  \item \textbf{Terza legge (periodi):} Il quadrato del periodo orbitale è
    proporzionale al cubo del semiasse maggiore:
    \[
      T^2=\frac{4\pi^2}{GM}\,a^3\,.
    \]
    Si ottiene integrando la seconda legge sull'area dell'ellisse
    $A=\pi ab$, usando le relazioni tra $a,b,L,E$.
\end{enumerate}

% ===============================================================
\section{Fluidi}
% ===============================================================

\subsection{Grandezze fondamentali dei fluidi}

Un fluido è descritto da: \textbf{densità} $\rho(\vec{r},t)$ [kg/m$^3$],
\textbf{pressione} $p(\vec{r},t)$ [Pa] (scalare, isotropa negli stati
d'equilibrio), \textbf{campo di velocità} $\vec{v}(\vec{r},t)$.

\subsection{Equazione della statica dei fluidi}

All'equilibrio la forza netta su ogni elemento di volume è nulla:
\[
  \nabla p=\rho\,\vec{g}\,.
\]
Per un fluido incompressibile ($\rho=\mathrm{const}$) con $\vec{g}=-g\hat{\bm{z}}$:
\[
  \frac{\mathrm{d}p}{\mathrm{d}z}=-\rho g\,.
\]

\subsection{Legge di Stevino e sua dimostrazione}

\begin{teorema}[Stevino]
In un fluido incompressibile in equilibrio, la differenza di pressione tra due
punti separati verticalmente da $h$ è:
\[
  \Delta p=\rho g h\,.
\]
\end{teorema}
\begin{proof}
Si isola un cilindretto di fluido di sezione $A$ e spessore $\mathrm{d}z$ a
quota $z$. Le forze che agiscono su di esso sono: pressione dal basso
($p(z)\,A$, verso l'alto), pressione dall'alto ($-p(z+\mathrm{d}z)\,A$),
peso ($-\rho g A\,\mathrm{d}z$). Equilibrio:
\[
  p(z)A-p(z+\mathrm{d}z)A-\rho g A\,\mathrm{d}z=0
  \implies \mathrm{d}p=-\rho g\,\mathrm{d}z\,.
\]
Integrando tra $z_1$ e $z_2=z_1+h$:
$p(z_1)-p(z_2)=\rho g h$, cioè $\Delta p=\rho g h$.
\end{proof}

\subsection{Legge di continuità}

Esprime la conservazione della massa:
\[
  \frac{\partial\rho}{\partial t}+\nabla\cdot(\rho\vec{v})=0\,.
\]
Per fluido incompressibile in regime stazionario ($\partial\rho/\partial t=0$,
$\rho=\mathrm{const}$): $\nabla\cdot\vec{v}=0$. Per un condotto con sezione
$A(x)$ variabile:
\[
  A_1 v_1=A_2 v_2=\mathrm{const}\,.
\]
La portata volumetrica $Q=Av$ si conserva: dove la sezione si restringe, la
velocità aumenta.

\subsection{Equazione di Bernoulli}

\begin{teorema}[Bernoulli]
Per un fluido ideale (incomprimibile, non viscoso), in moto stazionario, lungo
una linea di flusso:
\[
  p+\frac{1}{2}\rho v^2+\rho gz=\mathrm{const}\,.
\]
\end{teorema}
\textbf{Significato fisico:} i tre termini rappresentano rispettivamente
l'energia di pressione, l'energia cinetica e l'energia potenziale per unità
di volume. L'equazione è la conservazione dell'energia per unità di volume
lungo una streamline.

\textbf{Condizioni di validità:}
fluido incomprimibile, non viscoso (ideale), moto stazionario, applicata
lungo la stessa linea di flusso (non valida tra streamline diverse per fluidi
rotazionali).

% ===============================================================
\section{Onde e principio di sovrapposizione}
% ===============================================================

\subsection{Equazione di D'Alembert}

L'\textbf{equazione delle onde} in una dimensione è:
\[
  \frac{\partial^2 u}{\partial x^2}=\frac{1}{v^2}\frac{\partial^2 u}{\partial t^2}\,,
\]
dove $v$ è la velocità di propagazione del disturbo. È un'equazione alle
derivate parziali del secondo ordine, lineare e iperbolica.

\paragraph{Soluzione generale.}
La soluzione generale (D'Alembert, 1747) è:
\[
  u(x,t)=f(x-vt)+g(x+vt)\,,
\]
somma di un'onda progressiva $f(x-vt)$ (viaggia nella direzione $+x$) e di
un'onda regressiva $g(x+vt)$ (viaggia nella direzione $-x$), con $f$ e $g$
funzioni arbitrarie determinate dalle condizioni iniziali (profilo e velocità
del disturbo a $t=0$).

\subsection{Onde trasversali su una corda tesa}

Sia $\mathcal{T}$ la tensione e $\mu$ la densità lineare (massa/lunghezza)
della corda. Per un elemento di corda di lunghezza $\mathrm{d}x$, le
componenti trasversali delle tensioni alle due estremità danno una forza netta:
$\mathrm{d}F_y=\mathcal{T}\frac{\partial^2 y}{\partial x^2}\mathrm{d}x$.
Per il secondo principio: $\mu\,\mathrm{d}x\,\frac{\partial^2 y}{\partial t^2}=\mathcal{T}\frac{\partial^2 y}{\partial x^2}\mathrm{d}x$, cioè:
\[
  \frac{\partial^2 y}{\partial x^2}=\frac{\mu}{\mathcal{T}}\frac{\partial^2 y}{\partial t^2}\,,
\]
che è l'equazione di D'Alembert con \textbf{velocità di propagazione}:
\[
  v=\sqrt{\frac{\mathcal{T}}{\mu}}\,.
\]
Aumentando la tensione o riducendo la densità lineare, l'onda si propaga più
rapidamente.

\subsection{Principio di sovrapposizione}

\begin{teorema}[Sovrapposizione]
Se $u_1$ e $u_2$ sono soluzioni dell'equazione di D'Alembert, allora
$u=\alpha u_1+\beta u_2$ è anch'essa una soluzione, per ogni
$\alpha,\beta\in\mathbb{R}$.
\end{teorema}
\begin{proof}
L'equazione di D'Alembert è \emph{lineare omogenea}. Definendo l'operatore
$\mathcal{L}=\frac{\partial^2}{\partial x^2}-\frac{1}{v^2}\frac{\partial^2}{\partial t^2}$,
la proprietà di linearità dà:
$\mathcal{L}(\alpha u_1+\beta u_2)=\alpha\mathcal{L}u_1+\beta\mathcal{L}u_2=0$.
\end{proof}
\textbf{Significato fisico:} le onde si sommano algebricamente quando si
incontrano (interferenza costruttiva o distruttiva), continuano indisturbate
dopo l'incontro e non interagiscono tra loro. Il principio vale solo in regime
lineare: per onde di grande ampiezza (regime non lineare) la sovrapposizione
non è più valida e compaiono fenomeni come la formazione di shock.

\end{document}